% !Mode:: "TeX:UTF-8" 
\begin{conclusions}

本文围绕大语言模型在复杂问答生成过程中的可靠性问题，研究了过程错误识别、基于反思验证数据的错误缓解策略以及多模型协同可信问答流程。
与仅关注最终答案正确性的评测和优化方法不同，本文将研究重点放在模型的“思考-检索-信息生成-回答”过程，
关注模型如何调用自身参数化知识、错误如何在中间过程产生并传播，以及如何将错误分析结果转化为可用于训练和系统运行的监督信号。

第二章中，本文建立了面向自检索问答过程的错误识别方法。
本文基于 SSRL rollout 数据构造过程错误评测集，并比较了 PRM 异常检测和 LLM-as-a-judge 两类方法。
实验表明，基于 Qwen2.5-Math-PRM-7B 的 Z-score 异常检测具有较高精确率，但召回率较低，难以覆盖大量真实过程错误；
相比之下，基于评判模型的方法能够综合利用问题、完整生成轨迹和标准答案之间的语义关系。
说明反思验证数据能够有效提升小规模模型的过程级错误识别能力。

第三章中，本文提出基于反思验证数据的错误缓解策略。
本文使用强验证模型对自检索生成轨迹进行错误标注、错误分析和反思反馈生成，
构造 10000 条反思验证样本，并将其用于 Qwen2.5-3B-Instruct 的监督微调。
在此基础上，本文继续执行 SSRL 强化学习，使模型在结果奖励优化前先获得过程级错误意识。
实验结果表明，该策略在 PopQA、MuSiQue 和 BamboogleQA 等数据集上提升了最佳检查点表现，
这说明反思验证监督能够在一定程度上缓解结果奖励信号稀疏的问题，为后续强化学习提供更有利的初始化。

本文在第四章构建了多模型协同的迭代式可信问答流程。
本文将策略模型、标准答案匹配模块和验证模型结合起来，形成“生成-验证-反馈-再生成”的闭环。
策略模型负责生成多个候选自检索轨迹，验证模型负责检查推理过程、知识调用和最终答案之间的一致性，
未通过验证的样本则利用反馈进入下一轮反思式生成。实验结果表明，多模型协同验证能够在答案命中之外提供更严格的过程质量控制。

本文工作从错误识别、错误缓解到可信生成流程三个层面形成了较完整的研究链条：
首先利用过程级错误识别揭示模型生成轨迹中的可靠性风险；
其次将错误描述和反思反馈转化为训练信号，提升模型对自身知识调用错误的感知能力；
最后通过生成模型与验证模型的协同运行，在系统层面提高候选答案的可用性和可信度。
在未来，上述工作可以向大语言模型生成内容的评估、高质量数据构建和可信问答系统设计拓展并提供支持。



\end{conclusions}
